%%==========================================================================%%
%% background.tex -- "Background & Summary".                                 %%
%% Structure/boilerplate modeled on the reference Scientific Data paper.     %%
%% Prose paraphrased and genericized for whole-body pediatric PET/MRI.       %%
%% REWRITE and verify all claims/citations before submission. [TBD: ...] =   %%
%% fill in.                                                                   %%
%%==========================================================================%%

\section*{Background \& Summary}\label{sec:background}

Whole-body positron emission tomography combined with magnetic resonance imaging (PET/MRI)
pairs high-resolution, soft-tissue-rich anatomical information with a functional measure of metabolic activity.
The MRI component localises and characterises tissue without exposing patients to the ionising radiation of computed tomography (CT),
while the PET component quantifies tracer uptake and so distinguishes active disease from treated or inactive tissue.
These complementary properties make PET/MRI central to the diagnosis, staging, and treatment monitoring of paediatric malignancies such as lymphoma,
where avoiding cumulative radiation across repeated examinations is a primary concern [TBD: cite].

% Reusable framing: image-based AI -> upstream need for whole-body anatomical
% taxonomy + functional disease readout -> adult-only / CT-dominated gap ->
% pediatric PET/MRI need (mirrors the abstract argument). Expand here.
Image-based artificial intelligence (AI) tools are increasingly used across
clinical applications, and a key upstream capability for many of them is the
accurate identification of anatomy across the whole body, against which functional
findings such as metabolically active lesions can then be localised and scored.
Progress on these tasks has been enabled by several large imaging datasets.
To date, however, these have been overwhelmingly limited to adult subjects and dominated by anatomical CT,
with little paired functional imaging [TBD: cite].
This adult, CT-centric focus limits paediatric adoption for three reasons.
First, paediatric anatomy differs substantially from adult anatomy across development,
so models calibrated on adults alone behave unpredictably on paediatric inputs [TBD: cite].
Second, because of children's heightened sensitivity to ionising radiation,
MRI is preferred over CT for the anatomical component of paediatric whole-body imaging,
so adult CT corpora are of limited transfer value [TBD: cite].
Third, anatomy alone does not capture disease activity,
and the functional PET signal needed to detect lesions and assess treatment response is absent from these resources [TBD: cite].

% Contribution paragraph.
To help address these entangled gaps in image coverage, appropriate modality,
and functional disease labelling,
we present a paediatric whole-body PET/MRI dataset of 250 whole-body examinations from
134 patients, spanning ages 1 to 25 years (median 14).
The annotated subsets of this collection are being expanded, and the released cohort is
expected to grow as further annotations are completed.
The dataset includes both disease-free follow-up
examinations and examinations exhibiting active lymphoma, and is
distributed in NIfTI format together with [TBD: rich indexing metadata,
multi-reader anatomy segmentations, lesion annotations, staging and
metabolic-response labels, and (optionally) paired de-identified radiology
reports]. An overview of the cohort and the imaging protocols is shown in
Figure~\ref{fig:cohort}.

% Reuse / impact paragraph: what the data can be used for.
We anticipate that this dataset will support the development and benchmarking of
paediatric whole-body anatomical and oncological analysis tools, including [TBD: anatomical
segmentation, PET-guided lesion detection, automated staging and treatment-response
assessment, cross-modality transfer, and report-grounded multimodal modelling].
By providing paediatric PET/MRI with consistent organisation
and labelling, it complements existing adult, CT-dominated resources and lowers
the barrier to validating AI methods in paediatric settings.

%%--------------------------------------------------------------------------%%
%% Figure placeholder -- replace with our own figure (cf. Fig. 1 of ref:    %%
%% distribution of presentations/pathologies and imaging protocols).        %%
%%--------------------------------------------------------------------------%%
\begin{figure}[t]
\centering
\fbox{\parbox[c][4.5cm][c]{0.85\textwidth}{\centering\itshape
[Figure placeholder: cohort overview --- distribution of clinical
presentations/pathologies (a) and imaging protocols/sequences per subject (b).]}}
\caption{Overview of the cohort: distribution of clinical presentations and of
the imaging protocols/sequences represented in the dataset. [TBD: finalise.]}
\label{fig:cohort}
\end{figure}
